\documentclass[11pt]{article}
\usepackage[margin=1in]{geometry}
\usepackage{amsmath,amssymb,amsthm}
\usepackage[T1]{fontenc}
\usepackage{lmodern}
\usepackage[hidelinks]{hyperref}
\newtheorem{theorem}{Theorem}
\title{A scalar-matrix counterexample to Conjecture 00000002343}
\author{gaochengzhecpu}
\date{October 9, 2026}
\begin{document}
\maketitle
\begin{abstract}
For every random pencil of one-by-one matrices, the number of eigenvalue
collisions is identically zero, independently of the probability law.
Its expectation therefore differs from the conjectured value $\pi/4$.
The supplied Lean project verifies the actual characteristic polynomial,
absence of repeated roots, collision count, and expectation.
\end{abstract}

\section{Statement and scope}
The source conjecture asserts that the expected number of coincidences
(crossings) of eigenvalue curves of a random matrix pencil $A+\lambda B$
is $(\pi/4)n$. It imposes no lower dimension restriction beyond being a
matrix pencil, no asymptotic qualification, and no probability law.
We refute this asserted equality at the admissible dimension $n=1$.
The counterexample works for every law of $A$ and $B$, including
nondegenerate scalar Gaussian laws; it does not rely on choosing
deterministic random matrices.

At a parameter $\lambda$, a collision of distinct eigenvalue branches
requires a repeated eigenvalue. For a matrix $M$, this is a common zero
of its characteristic polynomial $p_M(z)=\det(zI-M)$ and its derivative.
We show that no such zero exists at any parameter. Consequently, whether
one counts collision parameters, colliding pairs, or transverse crossings,
the count in this example is zero.

\section{Counterexample}
\begin{theorem}
Let $(\Omega,\mathcal F,\mu)$ be any probability space and let
$a,b:\Omega\to\mathbb C$ be any scalar random variables.
For the one-by-one pencil
\[
 M_\omega(\lambda)=[a(\omega)+\lambda b(\omega)],\qquad \lambda\in\mathbb R,
\]
the collision set is empty for every $\omega$. The expected collision
count is $0$, not $\pi/4$.
\end{theorem}
\begin{proof}
Fix $\omega$ and $\lambda$. Direct computation of the one-by-one determinant
gives
\[
 p_{M_\omega(\lambda)}(z)=z-a(\omega)-\lambda b(\omega),
 \qquad p'_{M_\omega(\lambda)}(z)=1.
\]
Thus the unique eigenvalue is $a(\omega)+\lambda b(\omega)$, and it is simple.
Since the derivative is never zero, no common zero of $p$ and $p'$ exists.
The set of real parameters with a collision is therefore empty.
Its cardinality $N(\omega)$ is zero, and the random variable $N$ is the
measurable constant zero. Hence
\[
 \mathbb E_\mu[N]=\int_\Omega 0\,d\mu=0.
\]
The conjecture instead gives $(\pi/4)\cdot 1>0$, a contradiction.
\end{proof}

\section{Formal verification}
The file \texttt{lean/Main.lean} uses Mathlib's actual
\texttt{Matrix.charpoly} and the Bochner integral.
It proves the following chain, without a hypothesis asserting the desired
spectral conclusion:
\begin{itemize}
\item \texttt{pencil\_\allowbreak charpoly} evaluates the characteristic polynomial of
      arbitrary complex one-by-one matrices $A+\lambda B$.
\item \texttt{pencil\_\allowbreak unique\_\allowbreak eigenvalue} identifies its unique root;
      \texttt{pencil\_\allowbreak charpoly\_\allowbreak separable} proves separability.
\item \texttt{collision\_\allowbreak set\_\allowbreak empty} proves that the set defined by
      simultaneous vanishing of the characteristic polynomial and derivative
      is empty. Its \texttt{Nat.card} is consequently an ordinary finite
      cardinality; no convention for infinite cardinalities is used.
\item \texttt{expected\_\allowbreak crossing\_\allowbreak count\_\allowbreak zero} evaluates the actual integral
      for every measure and every pair of matrix-valued sample functions.
      \texttt{conjectured\_\allowbreak expectation\_\allowbreak false} gives the contradiction for
      every probability measure.
\end{itemize}
All final theorems use only Lean's standard logical axioms. The project
contains no unproved placeholders, custom axioms, or native evaluation
used as a proof. No auxiliary computation is required.

\paragraph{Reproduction.}
In \texttt{lean/}, run:
\begin{verbatim}
lake build
lake env lean -DwarningAsError=true Main.lean
\end{verbatim}
Lean 4.19.0 and the exact Mathlib revision are pinned in the project.
Run \texttt{tectonic main.tex} from the submission directory to rebuild
this PDF. The bilingual source is preserved in \texttt{SOURCE.md}.

\begin{thebibliography}{9}
\bibitem{source} The Last Math Competition,
\href{https://github.com/The-Last-Math-Competition/The-Last-Math-Competition/blob/main/conjectures/00000002343.md}{Conjecture 00000002343}.
\bibitem{mathlib} The Mathlib Community, Lean 4 mathematical library,
\href{https://github.com/leanprover-community/mathlib4/commit/c44e0c8ee63ca166450922a373c7409c5d26b00b}{revision c44e0c8ee63c}.
\end{thebibliography}
\end{document}
